\chapter{Alternative and Text Data}\label{ch:rs:alternative-and-text-data}

A vendor offers five years of card-spending data on four hundred retailers for a fee of several hundred thousand dollars a
year. The trial file arrives with a note: the panel changed card provider in its third year. The history looks excellent. The
questions are the chapter's: what exactly is this dataset, which of its history was known at the time, what does it add to
what the firm already has, is it legal to use, and what is it worth. Froot, Kang, Ozik and Sadka built real-time proxies of
retail sales from sources including about 50 million mobile devices and found that their within-quarter measure explained
sales growth and \omterm{def:rs:fundamental-analyst-and-event-features:sue}{earnings surprises}, with average excess returns at announcement of 3.4\%; Zhu (2019) found that the
introduction of consumer-transaction and satellite data increased the informativeness of prices. Such data work. The purchase is
a research project of its own, with a harness (\texttt{firm.vendoreval}), a synthetic trial whose truth is planted, and a decision
at the end.

\section{What the datasets are}

\begin{definition}[Alternative data, panel drift]\label{def:rs:alternative-and-text-data:altdata}
\emph{Alternative data}\index{alternative data} are data about companies, economies or markets from sources other than prices,
filings and official statistics: card and bank transactions, receipts, web traffic and prices, app usage, geolocation, satellite
and aerial images, shipping and flight records, job postings, text. A dataset built on a panel (of cardholders, devices, stores)
suffers \emph{panel drift}\index{panel drift} when the panel's composition or its relation to the quantity it measures changes over
time: a provider change, attrition, a new source, a re-weighting.
\end{definition}

Every alternative dataset is a measurement process with a history of its own, and three questions come before any backtest. What is
measured, by whom, and with what consent (the provenance)? When was each historical value first delivered (the \omterm{def:rs:point-in-time-data-and-the-biases:bitemporal}{knowledge time} of
chapter~3)? How has the panel changed (its drift)? The last two are usually tangled: a vendor launches, then reconstructs years of
history from today's panel to sell a longer backtest, and the reconstruction is fitted, knowingly or not, to outcomes that were known
when it was built. That is \omterm{def:rs:point-in-time-data-and-the-biases:vintage}{backfill bias} (chapter~3) in its commercial form.

\section{Evaluating a vendor}

\begin{definition}[Data trial, incremental information coefficient]\label{def:rs:alternative-and-text-data:trial}
A \emph{data trial}\index{data trial} is a time-limited evaluation of a dataset under a trial licence, before purchase. The
\emph{incremental information coefficient}\index{incremental information coefficient} of a new feature is the \omterm{def:rs:anatomy-of-a-predictor:ic}{information
coefficient} of its residual after regressing it, cross-section by cross-section, on the features the firm already owns.
\end{definition}

The chapter's trial is synthetic and its story planted, so every conclusion can be checked. In \texttt{firm.synthmkt}, 400 names
listed for all ten years are the retailers; each quarter's true \omterm{def:rs:fundamental-analyst-and-event-features:sue}{earnings surprise} is announced a random number of days after the
quarter's end, with a price jump in its direction. The vendor's panel value for each retailer and quarter is delivered ten trading days
after the quarter ends. Its five-year history (years 6 to 10 of the market) hides three things: the vendor went live when it changed
provider, at the start of the history's third year; the first two years were reconstructed at launch from the new provider's panel,
re-weighted to match the retailers' reported results; and the live panel covers fewer names, with a level shift and more noise. The firm
already owns a signal that carries part of the same information. The trial keeps only quarters whose announcement falls after the
delivery, 5\,980 rows, and \texttt{firm.vendoreval} measures what it can see.

\textbf{Coverage and history.} The backfilled years have 338 rows a quarter, the live years 273: the sample shrank when the product went
live, the opposite of a growing panel, and a clue. Of all rows, 39.6\% were first delivered more than a quarter after their period (the
vendor's own launch date, which the trial must ask for, dates them).

\textbf{The relation to the truth.} Regressing the panel value on the reported surprise year by year
(\cref{fig:rs:alternative-and-text-data:drift}): slope 0.92 and 0.87 in the first two years, 0.37, 0.39 and 0.35 after; correlation
0.85 and 0.84, then 0.35, 0.35 and 0.32; an intercept of 0.02 before and 0.12 to 0.17 after. A CUSUM of the residuals, ordered in time (chapter~13),
puts the single most likely break at the start of the third year, with a statistic of 3.05 (the 5\% critical value is about 1.36). The
provider change the vendor mentioned is also where the history stops being a reconstruction.

\begin{omfigure}
\begin{tikzpicture}
\begin{axis}[width=10cm,height=6cm,xlabel={reported earnings surprise (standardised)},ylabel={vendor's panel value},
  tick label style={font=\small},label style={font=\small},xmin=-3.5,xmax=3.5,ymin=-4,ymax=4.5,grid=major,grid style={gray!25},
  legend columns=3,
  legend style={font=\footnotesize,at={(0.5,-0.28)},anchor=north,draw=none,fill=none,
  /tikz/every even column/.append style={column sep=8pt}},table/col sep=comma]
\addplot[only marks,mark=*,mark size=0.7pt,gray!70] table[x=s,y=m]{figdata/research/12-alternative-and-text-data/live.csv};
\addplot[only marks,mark=*,mark size=0.7pt,omDef] table[x=s,y=m]{figdata/research/12-alternative-and-text-data/backfill.csv};
\addplot[omThm,very thick] table[x=s,y=m]{figdata/research/12-alternative-and-text-data/fitline.csv};
\legend{live years 3--5,backfilled years 1--2,live fit}
\end{axis}
\end{tikzpicture}

\omcaption{The vendor's panel value against the retailers' reported surprise: a sample of 600 rows from each regime. The backfilled
history hugs the outcome it was fitted to; the live panel is shifted and noisier. Data: the chapter's synthetic trial on
\texttt{firm.synthmkt}.}\label{fig:rs:alternative-and-text-data:drift}
\end{omfigure}

\textbf{The \omterm{def:rs:anatomy-of-a-predictor:ic}{information coefficient} by year} with the return from delivery to announcement: 0.38 and 0.37 in the backfilled years, 0.12,
0.14 and 0.15 in the live ones (\cref{fig:rs:alternative-and-text-data:ic}). A backtest over the whole history would report the average of
the two regimes and describe neither.

\textbf{What is new.} In the live years the firm's own signal has an IC of 0.20, larger than the vendor's 0.13. Residualised on it quarter by
quarter, the vendor's feature keeps an incremental IC of 0.10, with a $t$ statistic of 4.0 over 12 quarters: the panel carries information the
firm lacks, less than its raw IC suggests.

\begin{omfigure}
\begin{tikzpicture}
\begin{axis}[width=10cm,height=5.4cm,bar width=14pt,xlabel={year of the vendor's history},ylabel={rank IC},tick label style={font=\small},
  label style={font=\small},xtick={1,2,3,4,5},xmin=0.4,xmax=5.6,ymin=0,ymax=0.45,grid=major,grid style={gray!25},
  legend columns=3,
  legend style={font=\footnotesize,at={(0.5,-0.28)},anchor=north,draw=none,fill=none,
  /tikz/every even column/.append style={column sep=8pt}},table/col sep=comma]
\addplot[ybar,area legend,fill=omDef!55,draw=omDef] table[x=year,y=ic]{figdata/research/12-alternative-and-text-data/ic.csv};
\addplot[omThm,thick,dashed,sharp plot,mark=none] table[x=year,y=inc]{figdata/research/12-alternative-and-text-data/icref.csv};
\addplot[black!60,thick,dotted,sharp plot,mark=none] table[x=year,y=own]{figdata/research/12-alternative-and-text-data/icref.csv};
\legend{vendor's panel,incremental IC (live),firm's own signal (live)}
\end{axis}
\end{tikzpicture}

\omcaption{Rank IC of the vendor's panel value with the return from delivery to announcement, by year of its history; for the live years,
the IC of the signal the firm already owns and the vendor's incremental IC over it. Data: the chapter's synthetic
trial.}\label{fig:rs:alternative-and-text-data:ic}
\end{omfigure}

The real-data counterparts of each step are well known to anyone who has run trials: coverage tables with gaps where a merchant left the
panel, histories whose first delivery date the vendor cannot supply, tickers mapped with today's \omterm{def:rs:universe-symbology-and-corporate-actions:master}{security master} (chapter~4), data revised
without notice. The harness asks the same questions of any file.

\section{Text as data}

\begin{definition}[Dictionary method, sentiment score, document embedding]\label{def:rs:alternative-and-text-data:text}
A \emph{dictionary method}\index{dictionary method} scores a document by counting the words it contains from fixed word lists. A
\emph{sentiment score}\index{sentiment score} is such a measure of tone, for instance the share of negative words, or negative minus
positive words over all words. A \emph{document embedding}\index{document embedding} maps a document to a vector of real numbers (word
counts weighted by rarity, or the output of a trained language model) so that documents with similar content have nearby vectors.
\end{definition}

Tetlock (2007) measured the pessimism of a daily \emph{Wall Street Journal} column and found that high media pessimism predicts downward
pressure on market prices followed by a reversion to fundamentals, and that unusually high or low pessimism predicts high trading volume.
Loughran and McDonald (2011) showed that word lists built for other disciplines misclassify common words in financial text: in 10-Ks from
1994 to 2008, almost three-fourths of the words the widely used Harvard dictionary classifies as negative are not negative in financial
contexts (liability, tax, cost, capital), and they built finance-specific lists.

The chapter's corpus plants that problem. Two thousand synthetic documents of 400 words have a tone $\tau$ that raises the rate of eight
positive words and lowers that of eight negative ones; eight finance words that a general list calls negative are frequent and unrelated to
tone, and make 77\% of the general list's negative hits. The negative-word share correlates with $-\tau$ at 0.74 with the general list, 0.87
with the finance list, and 0.95 for negative minus positive words with the finance lists. The numbers are the corpus's; the lesson is Loughran
and McDonald's: a dictionary is a measurement instrument, calibrated for a domain.

Embeddings replace hand-built lists with learned representations and move the problems rather than solving them. Their training corpus has a
date, and a model trained on text written after the backtest's decision date knows how the story ended; its vocabulary and its sense of what
is similar drift with each version; and the mapping from embedding to return is one more fitted model, with the overfitting risks of
chapter~20. The point-in-time rule applies to the model as to the data: the model used at a decision date is one that existed then.

\section{Law and ethics}

\begin{definition}[Material non-public information]\label{def:rs:alternative-and-text-data:mnpi}
\emph{Material non-public information}\index{material non-public information} (MNPI) is information about a security or its issuer that has
not been made public and that is material: in the SEC's words, a matter is material if there is a substantial likelihood that a reasonable
person would consider it important. Trading on MNPI obtained in breach of a duty is insider trading.
\end{definition}

\omterm{def:rs:alternative-and-text-data:altdata}{Alternative data} sit near that line by design: they are valuable because they are not in the filings. The legal questions are about how the
data were obtained and whether the seller had the right to sell them, and a buyer's diligence is part of its own defence.

\begin{dated}{2026-09}{Alternative data in enforcement and litigation}\label{dat:rs:alternative-and-text-data:law}
In September 2021 the SEC settled securities fraud charges against App Annie, an \omterm{def:rs:alternative-and-text-data:altdata}{alternative data} provider for the mobile app industry, and
its co-founder, who agreed to pay more than \$10 million: the SEC's first enforcement action charging an \omterm{def:rs:alternative-and-text-data:altdata}{alternative data} provider with
securities fraud. The order found that, while telling companies their confidential app data would be aggregated and anonymised, App Annie used
non-aggregated, non-anonymised data from late 2014 to mid-2018 to alter its estimates and make them more valuable to trading firms, and
misrepresented this to those firms. In April 2022 the US Court of Appeals for the Ninth Circuit, in \emph{hiQ Labs v.\ LinkedIn}, again
affirmed a preliminary injunction letting a data analytics company access publicly available LinkedIn profiles, holding that the Computer
Fraud and Abuse Act's concept of access ``without authorization'' does not apply to public websites. In the European Union, the General Data
Protection Regulation (Regulation 2016/679) defines personal data as any information relating to an identified or identifiable natural
person, identifiable directly or indirectly, including by location data or an online identifier.
\end{dated}

A research team's checklist follows from these cases: the vendor's written representation of how the data were collected and with which
consents; whether any source is an insider or bound by confidentiality; whether personal data are present, pseudonymised or truly aggregated;
the terms of service of scraped sites (a preliminary injunction on public profiles under one statute is not a licence to scrape); a legal and
compliance sign-off before the trial, not after the purchase; and a record of all of it, kept with the dataset's metadata. The ethics are wider
than the law: data that let a desk infer who visited a clinic are not acceptable because a court has not yet ruled on them.

\section{The buy decision}

The value of a dataset is the profit of what it adds, net of costs, over the life of the contract, as the edge decays because other buyers
trade on it. In the trial, the long--short spread of the top and bottom quintiles of the incremental feature earns 2.9\% a quarter in the
live years (live raw feature: 3.4\%; the backfilled years: 10.4\%). With the desk's capacity for announcement trades in these names put at
\$20 million, trading costs of 0.4\% a year and a two-year half-life for the edge, the break-even annual fee over a three-year contract is
\$1.35 million. The vendor's history, taken at face value, implies \$5.06 million; the live raw feature \$1.59 million
(\cref{fig:rs:alternative-and-text-data:breakeven}). The half-life decides the rest: at six months the incremental break-even is \$0.43
million, near the asking price.

\begin{omfigure}
\begin{tikzpicture}
\begin{axis}[width=10cm,height=5.6cm,xlabel={half-life of the edge (years)},ylabel={break-even fee (\$ million a year)},
  tick label style={font=\small},label style={font=\small},xmin=0.4,xmax=5.1,ymin=0,ymax=7,grid=major,grid style={gray!25},
  legend columns=3,
  legend style={font=\footnotesize,at={(0.5,-0.28)},anchor=north,draw=none,fill=none,
  /tikz/every even column/.append style={column sep=8pt}},table/col sep=comma]
\addplot[gray,thick,dashed,mark=o] table[x=half_life,y=back]{figdata/research/12-alternative-and-text-data/breakeven.csv};
\addplot[omDef,thick,mark=*] table[x=half_life,y=live]{figdata/research/12-alternative-and-text-data/breakeven.csv};
\addplot[omThm,very thick,mark=square*] table[x=half_life,y=inc]{figdata/research/12-alternative-and-text-data/breakeven.csv};
\legend{backfilled history (the pitch),live years,live and incremental}
\end{axis}
\end{tikzpicture}

\omcaption{Break-even annual fee of the synthetic card panel over a three-year contract, against the half-life of its edge, for \$20 million
of capacity and 0.4\% a year of trading costs: priced on the vendor's backfilled history, on the live years, and on the live years net of the
firm's own signal. Data: the chapter's synthetic trial.}\label{fig:rs:alternative-and-text-data:breakeven}
\end{omfigure}

Three corrections separate the pitch from the price: live data only (the backfill tripled the spread), incremental over what is owned (a
further 15\%), and decay (a factor of three between a six-month and a two-year half-life). The half-life is the hardest to know before
buying: it depends on how many others buy the same file, which the vendor knows and the buyer does not; chapter~28 measures it after the fact.
Capacity (chapter~28) and costs (chapter~27) come from the firm's own models. The chapter's weekend problem makes the decision.

\section{Predictor cards}

\begin{predictorcard}{Card-panel sales surprise}\label{pred:rs:alternative-and-text-data:card}
\sfield{Definition} The panel's spending growth for the quarter, residualised on the firm's existing signals, ranked across the retailers
announcing in the next weeks.
\sfield{Inputs and timestamps} The vendor's quarterly values with their first delivery date (never the period); the announcement calendar
(chapter~11).
\sfield{Rationale} Within-quarter sales proxies explain surprises and announcement returns (Froot, Kang, Ozik and Sadka).
\sfield{Horizon and half-life} From delivery to the announcement; the edge's half-life is set by the number of buyers.
\sfield{Normalisation} Within the quarter's cross-section; a panel-change indicator as a control.
\sfield{Failure modes} Backfilled history; \omterm{def:rs:alternative-and-text-data:altdata}{panel drift}; coverage bias toward the panel's merchants; buyers crowding the same file.
\sfield{Sources} As cited; \texttt{rs\_altdata.trial\_report}.
\end{predictorcard}

\begin{predictorcard}{Finance-dictionary tone}\label{pred:rs:alternative-and-text-data:tone}
\sfield{Definition} Negative minus positive words from finance-specific lists, over all words, in a filing, release or news article.
\sfield{Inputs and timestamps} The document's publication time (the filing's acceptance time, the article's time stamp); the word lists'
version.
\sfield{Rationale} Tone carries information or moves prices through sentiment (Tetlock; Loughran and McDonald).
\sfield{Horizon and half-life} Days after the document; long reversals for sentiment-driven moves.
\sfield{Normalisation} Relative to the same firm's past documents of the same type (boilerplate is stable).
\sfield{Failure modes} General-purpose lists; word lists or models built after the decision date; boilerplate changes.
\sfield{Sources} As cited; \texttt{rs\_altdata.text\_scores}.
\end{predictorcard}

\section{Tutorial: running a vendor trial}\label{tut:rs:alternative-and-text-data:tut}

\begin{tutorial}
\textbf{Goal.} Run the trial harness on the synthetic card panel: coverage, backfill share, the panel's relation to the truth by year and
its break, IC by year, incremental IC, long--short spreads and the break-even fee; then score the synthetic corpus with two word lists.
\textbf{End state:}
\cref{fig:rs:alternative-and-text-data:drift,fig:rs:alternative-and-text-data:ic,fig:rs:alternative-and-text-data:breakeven}; the text
correlations.
\begin{tutsteps}
\item \textbf{Locate the break}: a CUSUM of the residuals of the panel on the truth, in time order.
\omcode{code/firm/vendoreval/firm_vendoreval.py}{47}{53}{The most likely single break in the mean.}\label{lst:rs:alternative-and-text-data:cusum}
\item \textbf{Incremental IC}: residualise within each cross-section, then rank-correlate.
\omcode{code/firm/vendoreval/firm_vendoreval.py}{69}{83}{Residualising on owned signals, and the incremental IC.}\label{lst:rs:alternative-and-text-data:inc}
\item \textbf{Price it}: the fee that equals the mean annual net profit as the edge decays.
\omcode{code/firm/vendoreval/firm_vendoreval.py}{96}{101}{The break-even annual fee.}\label{lst:rs:alternative-and-text-data:fee}
\item \textbf{Run} \texttt{trial\_report()}, \texttt{breakeven\_curve()}, \texttt{text\_scores()} and \texttt{fig\_altdata.py}.
\end{tutsteps}
\textbf{What to change next.} Hide the launch date and try to find it from the data alone; give the firm's own signal more of the panel's
information and watch the incremental IC and the price fall.
\end{tutorial}

\section{Build: the vendor-trial harness}\label{bld:rs:alternative-and-text-data:vendoreval}

\begin{build}
\bfield{Purpose} The same questions asked of every dataset the firm considers, with answers that can be compared across vendors and filed with
the purchase decision.
\bfield{Interface} \texttt{coverage}, \texttt{backfill\_share(period, delivered, lag)}, \texttt{fit\_by\_group}, \texttt{cusum\_break},
\texttt{rank\_ic}, \texttt{ic\_by\_group}, \texttt{residualise}, \texttt{incremental\_ic(f, X, y, group)}, \texttt{long\_short},
\texttt{breakeven(annual, capital, cost, half\_life, years)}, \texttt{tone(counts, vocab, words)}.
\bfield{Rules} No trial result without first-delivery dates or an explicit statement that they are missing; ICs by period, never only pooled;
value computed on the incremental feature, net of costs and decay.
\bfield{Acceptance tests} \texttt{code/firm/vendoreval/tests/}: coverage and backfill share by hand; per-group fits; a planted mean shift found by
the CUSUM and none in noise; a residual uncorrelated with the owned signal, an incremental IC below the raw one and a copy with nothing new; the
fee without decay equals the annual profit; tone shares by hand.
\bfield{Stretch} Structural-break tests with several breaks; a Bayesian update of the half-life from the first months of live trading; a
provenance checklist stored with each dataset.
\end{build}

\begin{omsources}
\item K.~Froot, N.~Kang, G.~Ozik and R.~Sadka, ``What do measures of real-time corporate sales say about earnings surprises and post-announcement
returns?'', \emph{Journal of Financial Economics} 125(1), 2017.
\item C.~Zhu, ``Big data as a governance mechanism'', \emph{Review of Financial Studies} 32(5), 2019.
\item P.~C.~Tetlock, ``Giving content to investor sentiment: the role of media in the stock market'', \emph{Journal of Finance} 62(3), 2007.
\item T.~Loughran and B.~McDonald, ``When is a liability not a liability? Textual analysis, dictionaries, and 10-Ks'', \emph{Journal of Finance}
66(1), 2011.
\item US Securities and Exchange Commission, press release 2021-176 (App Annie), 14 September 2021; Staff Accounting Bulletin No.~99,
Materiality.
\item \emph{hiQ Labs, Inc.\ v.\ LinkedIn Corp.}, No.~17-16783 (9th Cir., 18 April 2022).
\item Regulation (EU) 2016/679 (General Data Protection Regulation), Article~4.
\end{omsources}

\section{Exercises}

\begin{exercise}[$\star$]\label{exo:rs:alternative-and-text-data:1}
A trial file has 1\,200 rows for periods before the vendor's launch and 1\,800 after. With no information on delivery dates, what share of the
history is at risk of backfill, and what would you ask the vendor for?
\end{exercise}

\begin{exercise}[$\star$]\label{exo:rs:alternative-and-text-data:2}
A document of 400 words contains 12 words from a general negative list, 9 of which are liability, tax, cost or capital, and 5 positive words.
Compute its negative share with the general list and with a list that excludes those four words, and its net tone with the second.
\end{exercise}

\begin{exercise}[$\star$]\label{exo:rs:alternative-and-text-data:3}
A strategy earns 2.9\% a quarter on \$20 million with 0.4\% a year of costs and no decay. What is the break-even annual fee?
\end{exercise}

\begin{exercise}[$\star\star$]\label{exo:rs:alternative-and-text-data:4}
The firm's own signal and the vendor's panel both measure the surprise $s$ with independent noise: $o = 0.45s + e_1$, $m = 0.35s + e_2$. Why is
the vendor's incremental IC positive but below its raw IC? What happens if $e_2$ is a copy of $e_1$?
\end{exercise}

\begin{exercise}[$\star\star$]\label{exo:rs:alternative-and-text-data:5}
Show that the break-even fee of a contract of $Y$ years, annual profit $A$ in the first year halving every $h$ years (measured at mid-year), and no
costs, is $\frac{A}{Y}\sum_{k=0}^{Y-1} 2^{-(k+1/2)/h}$. Evaluate it for $A = \$2.3$ million, $Y = 3$, $h = 2$.
\end{exercise}

\begin{exercise}[$\star\star$]\label{exo:rs:alternative-and-text-data:6}
Why does a \omterm{def:rs:alternative-and-text-data:text}{document embedding} from a language model trained in 2025 leak into a backtest of 2015 to 2020, even if every document is dated
correctly?
\end{exercise}

\begin{exercise}[$\star\star\star$]\label{exo:rs:alternative-and-text-data:7}
\emph{Coding.} In the synthetic trial, hide the flag of backfilled rows and locate the break from the panel's IC by quarter instead of its fit to
the truth. How sharp is the break, and why is the fit to the truth the better instrument?
\end{exercise}

\begin{exercise}[$\star\star\star$]\label{exo:rs:alternative-and-text-data:8}
\emph{Find the flaw.} ``The vendor confirmed the data are fully anonymised, and the source is a public website, so there is no legal question.''
\end{exercise}

\section{Problem: Should We Buy the Card Panel?}

\begin{problem}[{Weekend problem --- a vendor trial, priced}]\label{pb:rs:alternative-and-text-data:1}
The chapter's synthetic trial: 400 retailers, five years of a card panel, a provider change at the start of the third year, a signal the firm
already owns, an asking price of several hundred thousand dollars a year.

\medskip\noindent\textbf{Part I --- The file.}
\begin{enumerate}
\item How many rows does the trial keep, and why only quarters with an announcement after the delivery?
\item What are the rows per quarter before and after the launch, and why is the fall a clue?
\item What share of rows was delivered more than a quarter after its period?
\item What must the vendor provide for that share to be computed?
\end{enumerate}

\medskip\noindent\textbf{Part II --- The drift.}
\begin{enumerate}[resume]
\item Give the slope, correlation and intercept of the panel on the truth in each year.
\item Where does the CUSUM put the break, and with what statistic?
\item Why does the break coincide with the provider change here, and what else changed there?
\item What would you conclude if the vendor had not mentioned the provider change?
\end{enumerate}

\medskip\noindent\textbf{Part III --- The information.}
\begin{enumerate}[resume]
\item What are the ICs by year, and the IC of the firm's own signal in the live years?
\item What is the incremental IC and its $t$ statistic?
\item What are the long--short spreads of the backfilled years, the live raw feature and the incremental feature?
\item Why must the value use the incremental feature on live data only?
\end{enumerate}

\medskip\noindent\textbf{Part IV --- The decision.}
\begin{enumerate}[resume]
\item What break-even fees do the three spreads imply at a two-year half-life?
\item What is the incremental break-even at a half-life of six months and of five years?
\item State the \emph{named result}: the break-even annual price of the dataset given its incremental IC, capacity and decay.
\item Should the firm buy at \$400\,000 a year? At \$900\,000?
\item What would you negotiate for, besides the price?
\item Which legal and compliance questions must be answered before the trial?
\item How would you check the half-life after buying?
\item In one sentence: what is a dataset worth?
\end{enumerate}
\end{problem}

\section{Interview questions}

\begin{interviewq}[$\star$ \iqroles{researcher}]\label{iq:rs:alternative-and-text-data:1}
Name three ways an alternative dataset's history can mislead a backtest.
\end{interviewq}

\begin{interviewq}[$\star\star$ \iqroles{researcher}]\label{iq:rs:alternative-and-text-data:2}
How would you decide whether a new dataset adds anything to the signals you already have?
\end{interviewq}

\begin{interviewq}[$\star\star$ \iqroles{researcher, mle}]\label{iq:rs:alternative-and-text-data:3}
You score news with a language model. What point-in-time problems do you have to solve?
\end{interviewq}

\begin{interviewq}[$\star\star$ \iqroles{researcher, trader}]\label{iq:rs:alternative-and-text-data:4}
A vendor quotes \$500\,000 a year. How do you decide whether to pay?
\end{interviewq}

\begin{interviewq}[$\star\star$ \iqroles{researcher, risk}]\label{iq:rs:alternative-and-text-data:5}
What makes \omterm{def:rs:alternative-and-text-data:altdata}{alternative data} a legal risk, and what diligence would you do?
\end{interviewq}

\begin{interviewq}[$\star\star\star$ \iqroles{researcher}]\label{iq:rs:alternative-and-text-data:6}
Why do general-purpose sentiment dictionaries fail on financial text, and how would you build a better measure?
\end{interviewq}
